% ============================================================
% Section 142: 光激发过剩载流子分布与光电导
% ============================================================
\section{半导体物理：光激发过剩载流子分布与光电导}
\label{sec:photoconductivity}

\begin{modelbox}[题目：表面光激发半导体的载流子分布与光电导]
如图，用 $4880\,\text{\AA}$ Ar$^+$ 激光照射一块长度及宽度远大于厚度 $t$ 及扩散长度 $L$ 的半导体 Si。入射光在表面层内被完全吸收，样品上下表面复合速度分别为 $S_1$、$S_2$，单位面积光生载流子激发率为 $g$。

(1) 写出光激发产生过剩载流子运动方程及边界条件；

(2) 试证光生载流子沿厚度 $x$ 的分布为
\[
\Delta p(x)=\frac{Lg}{D}F_1,\quad F_1=\frac{\cosh\dfrac{t-x}{L}+a_2\sinh\dfrac{t-x}{L}}{(a_1+a_2)\cosh\dfrac{t}{L}+(1+a_1a_2)\sinh\dfrac{t}{L}},
\]
其中 $L$、$\tau$ 为扩散长度及寿命，$D$ 为双极扩散系数，$a_1=S_1L/D$，$a_2=S_2L/D$；

(3) 试求单位面积样品光电导
\[
\Delta G=eg\tau\mu_p(1+b)F_2,
\]
其中 $b=\mu_n/\mu_p$。
\end{modelbox}

\begin{tipbox}[考点快速分析]
\begin{itemize}
    \item \textbf{稳态双极扩散方程}：$D\dfrac{d^2\Delta p}{dx^2}-\dfrac{\Delta p}{\tau}=0$
    \item \textbf{边界条件}：上表面 $-D\dfrac{d\Delta p}{dx}\big|_0+S_1\Delta p(0)=g$；下表面 $-D\dfrac{d\Delta p}{dx}\big|_t-S_2\Delta p(t)=0$
    \item \textbf{双曲函数化简}：通解 $\Delta p(x)=Ae^{-x/L}+Be^{x/L}$
    \item \textbf{光电导}：$\Delta G=e(\mu_n\Delta n+\mu_p\Delta p)=e\mu_p(1+b)\Delta P$
\end{itemize}
\end{tipbox}

\begin{derivationbox}[标准解题过程]
\textbf{(1) 运动方程与边界条件}

过剩载流子浓度 $\Delta p(x)\approx\Delta n(x)$ 满足一维稳态双极扩散方程：
\[
D\frac{d^2\Delta p}{dx^2}-\frac{\Delta p}{\tau}=0,
\]
其中 $L=\sqrt{D\tau}$ 为扩散长度。

边界条件：
\begin{itemize}
    \item 上表面（$x=0$）：$-D\dfrac{d\Delta p}{dx}\big|_{x=0}+S_1\Delta p(0)=g$
    \item 下表面（$x=t$）：$-D\dfrac{d\Delta p}{dx}\big|_{x=t}-S_2\Delta p(t)=0$
\end{itemize}

\textbf{(2) 载流子分布的推导}

通解 $\Delta p(x)=Ae^{-x/L}+Be^{x/L}$。代入边界条件，定义 $a_1=S_1L/D$，$a_2=S_2L/D$，解出 $A$、$B$ 并化简得
\[
\Delta p(x)=\frac{Lg}{D}\cdot\frac{\cosh\dfrac{t-x}{L}+a_2\sinh\dfrac{t-x}{L}}{(a_1+a_2)\cosh\dfrac{t}{L}+(1+a_1a_2)\sinh\dfrac{t}{L}}=\frac{Lg}{D}F_1.
\]

\textbf{(3) 单位面积光电导}

单位面积内总的过剩空穴数：
\[
\Delta P=\int_0^t\Delta p(x)\,dx=\frac{Lg}{D}\cdot\frac{L\left(\sinh\dfrac{t}{L}+a_2\cosh\dfrac{t}{L}-a_2\right)}{(a_1+a_2)\cosh\dfrac{t}{L}+(1+a_1a_2)\sinh\dfrac{t}{L}}.
\]
单位面积光电导：
\[
\Delta G=e(\mu_n\Delta n+\mu_p\Delta p)=e\mu_p(1+b)\Delta P,
\]
其中 $b=\mu_n/\mu_p$，$\Delta n=\Delta p$。利用 $L^2=D\tau$ 化简得
\begin{equation}
\boxed{\Delta G=eg\tau\mu_p(1+b)F_2,}
\end{equation}
其中
\[
F_2=\frac{\sinh\dfrac{t}{L}+a_2\cosh\dfrac{t}{L}-a_2}{(a_1+a_2)\cosh\dfrac{t}{L}+(1+a_1a_2)\sinh\dfrac{t}{L}}.
\]
\end{derivationbox}

\begin{formulabox}[答案汇总]
\begin{center}
\begin{tabular}{ll}
\toprule
\textbf{项目} & \textbf{表达式} \\
\midrule
过剩载流子分布 & $\Delta p(x)=\dfrac{Lg}{D}F_1$ \\
单位面积光电导 & $\Delta G=eg\tau\mu_p(1+b)F_2$ \\
\bottomrule
\end{tabular}
\end{center}
\end{formulabox}

\begin{insightbox}[物理图像]
\begin{itemize}
    \item 表面复合速度 $S_1$、$S_2$ 对载流子分布和光电导有重要影响。无量纲参数 $a_1$、$a_2$ 表征表面复合与体扩散的相对强弱。
    \item 当样品厚度 $t\ll L$（薄样品）或 $t\gg L$（厚样品）时，$F_1$、$F_2$ 可进一步简化，对应不同的光电导响应。
    \item 双极扩散系数 $D$ 和双极迁移率是描述电子-空穴耦合输运的关键参数。本题结果在半导体光电器件（如光电导探测器、太阳电池）的建模中有广泛应用。
\end{itemize}
\end{insightbox}
